\noindent\textit{2008 manuscript.}

\section{Chemical Potential and Gibbs Distribution}

\subsubsection*{Chemical Potential and Entropy}

Let $\sigma = \sigma(U,V,N)$ (so now $\Sigma$ has 3 coordinates $U,V,N$).  Then
\[
d\sigma = \left( \frac{ \partial \sigma}{\partial U}\right)_{V,N}dU + \left( \frac{ \partial \sigma}{\partial V}\right)_{U,N} dV + \left( \frac{ \partial \sigma}{\partial N}\right)_{U,V}dN
\]

Consider curve $\begin{aligned} & \quad \\
  & c : \mathbb{R} \to \Sigma \\
  & c(t) \in \Sigma \end{aligned}$ s.t. $\dot{c} = \dot{U} \frac{ \partial }{ \partial U} + \dot{N} \frac{ \partial }{ \partial N} \in \mathfrak{X}(\Sigma)$ (i.e. $dV(\dot{c}) = 0$)

Let $\begin{aligned} & \quad \\
  & \tau = \tau(U,V, N) \in C^{\infty}(\Sigma) \\
  & d\tau \in \Omega^1(\Sigma) \end{aligned}$

Suppose $d\tau(\dot{c})=0$ (i.e. $c$ is a null curve of $d\tau$ or $c$ is an isothermal curve, or i.e. overall temperature change $d\tau $ is $0$)


\[
d\sigma(\dot{c}) = \left( \frac{ \partial \sigma }{ \partial U } \right) dU(\dot{c}) + \left( \frac{ \partial \sigma }{ \partial N} \right) dN(\dot{c}) \text{ or } \frac{d\sigma(\dot{c})}{ dN(\dot{c})} = \left( \frac{ \partial \sigma}{ \partial U} \right) \frac{dU(\dot{c})}{ dN(\dot{c})} + \left( \frac{ \partial \sigma}{ \partial N} \right)
\]

Now recall that originally, considering a system with only thermal contact to a huge heat reservoir/bath (it can be the universe), then only $N$,$U$ mattered in describing this system, and so $\sigma = \sigma(N,U)$, so that $\frac{1}{\tau} = \left( \frac{ \partial \sigma}{ \partial U}\right)_N$.  

Now \[
\frac{dU(\dot{c})}{dN(\dot{c})} = \frac{\dot{c}(U)}{\dot{c}(N)} = \frac{\dot{U}}{\dot{N}} = \left( \frac{ \partial U}{ \partial N} \right)_{\tau,V} 
\]
Similarly, 
\[
\frac{d\sigma(\dot{c})}{dN(\dot{c})} = \frac{\dot{c}(\sigma)}{ \dot{N}} = \left( \frac{\partial \sigma}{\partial N} \right)_{\tau,V}
\]
So
\[
\left( \frac{ \partial \sigma}{ \partial N} \right)_{\tau,V} = \frac{1}{\tau} \left( \frac{ \partial U}{ \partial N} \right)_{\tau,V} + \left( \frac{ \partial \sigma}{ \partial N} \right)_{U,V}
\]

Consider $F \equiv U-\tau \sigma$.  Now consider $dF = dU- \sigma d\tau - \tau d\sigma$.  Apply $dF$ onto $\dot{c}$, our vector field generated by our particular isothermal, and constant volume curve in $\Sigma$:
\[
dF(\dot{c}) = dU(\dot{c}) - \sigma d\tau (\dot{c}) - \tau d\sigma(\dot{c}) = \dot{U}- 0 - \tau d\sigma(\dot{c})
\]
The right hand side is 
\[
dF(\dot{c}) = \dot{U} \frac{ \partial F}{ \partial U} + \dot{N} \left( \frac{ \partial F}{ \partial N} \right)_{\tau,V} = \dot{U}(1) + \dot{N} \left( \frac{ \partial F}{ \partial N} \right)_{\tau,N}
\]
Left hand side is 
\[
\dot{U}- \tau \left( \left( \frac{ \partial \sigma}{ \partial N} \right)_{\tau,V} \dot{N} - \frac{1}{\tau} \left( \frac{ \partial U}{ \partial N} \right)_{\tau,V} \dot{N} + \left( \frac{ \partial \sigma}{ \partial U} \right)_{V,N} \dot{U} \right)
\]
Subtract $\dot{U}$ from both sides.  $\dot{U}$ was arbitrary (and so is $\dot{N}$), and so
\[
\Longrightarrow \mu \equiv \left( \frac{ \partial F}{ \partial N} \right)_{\tau,V} \equiv \left( \frac{\partial U}{ \partial N}\right)_{\tau,V} - \tau \left( \frac{ \partial \sigma}{\partial N}\right)_{\tau,V} \text{ or } \mu = -\tau \left( \frac{ \partial \sigma}{\partial N}\right)_{U,V}
\]

\subsubsection*{Thermodynamic identity} Generalize thermodynamic identity for $\sigma = \sigma(U,V,N)$
\[
\begin{aligned}
  d\sigma & = \left( \frac{ \partial \sigma}{\partial U} \right)_{V,N} dU + \left( \frac{ \partial \sigma}{ \partial V} \right)_{U,N} dV + \left( \frac{ \partial \sigma}{ \partial N} \right)_{U,V} dN = \\
  & = \frac{1}{\tau} dU + \frac{p}{\tau} dV - \frac{\mu}{\tau} dN \text{ or } dU = \tau d\sigma - p dV + \mu dN
\end{aligned}
\]

\solutionhead{1} \textbf{Centrifuge}.  

$T = \frac{1}{2} m \dot{r}^2 + \frac{1}{2} m r^2 \omega^2$.

Consider $\mu_{ext} = - \frac{1}{2} M r^2 \omega^2$  (negative so for bigger $r$ away from $r=0$ axis, lower chemical potential $\mu$, so to show ``centrifugal force'' outwards.  

$\mu_{tot} = \tau \ln{ (n/n_Q)} - \frac{1}{2} M r^2 \omega^2$.  

$\mu_{tot}(r) = \mu_{tot}(0)$ \emph{for diffusion equilibrium.}

\[
\tau \ln{\left( \frac{ n }{n_Q} \right)} - \frac{1}{2} M r^2 \omega^2 = \tau \ln{ \left( \frac{n(0)}{n_Q} \right)} \Longrightarrow \tau \ln{ \left( \frac{n(r)}{n(0)} \right) } =  \frac{1}{2} M r^2 \omega^2
\]

\[
\frac{n(r)}{n(0)} = \exp{ \left( \frac{Mr^2 \omega^2}{ 2\tau} \right) } \text{ or } \boxed{ n(r) = n(0) \exp{ \left( \frac{Mr^2 \omega^2}{2\tau} \right) } }
\]



\solutionhead{2} \textbf{Molecules in the Earth's atmosphere}.

Recall that for ideal gas, $F = - \tau [ N \ln{Z_1} - \ln{N!} ]$; $Z_1 = n_Q V = \left( \frac{ M \tau}{2\pi \hbar^2} \right)^{3/2} V$.  
\[
\mu = \left( \frac{ \partial F}{ \partial N} \right)_{\tau, V} = -\tau [ \log{Z_1} - \frac{d}{dN} \ln{N!} ] = \tau \ln{ \left( \frac{n}{n_Q} \right) }
\]
where $n_Q = \left( \frac{ M \tau}{2\pi \hbar^2} \right)^{3/2}$.

Now
$\mu_{\int} = \tau \log{ \left( \frac{n}{n_Q} \right)}$ \\
$\mu_{ext} = \frac{ - GM_e}{r} =  - \frac{ gR^2_m}{ r}$ since $g = \frac{ GM_e}{R_e^2}$.  
\[
\mu_{tot} = \tau \ln{ (n/n_Q)} + Mgh
\]

\emph{In equilibrium, this must be independent of $r$}: $\mu_{tot}(r) = \mu_{tot}(R)$.  

\[
\tau \ln{ ( n(r)/nQ)} - \frac{ gR^2 M }{r} = \tau \ln{ ( n(R)/n_Q) } - \frac{ MgR^2}{R}
\]
\[
\begin{gathered}
  \tau \ln{ \left( \frac{n(r)}{ n(R)} \right) } = Mg \left( \frac{R^2}{r} - R \right) \text{ or } \exp{ \left( \frac{Mg}{\tau} \left( \frac{R^2}{r} - R \right) \right) } = \frac{n(r)}{n(R)} \\ 
  \Longrightarrow n(r) = n(R) \exp{ \left( \frac{Mg}{ \tau} \left( \frac{R^2}{r} - \frac{1}{R} \right)\right)}
\end{gathered}
\]
so that 
\[
\boxed{ N = 4\pi n(R) \exp{ \frac{ - MgR}{\tau} \int_R^{\infty} r^2 dr \exp{ \left( \frac{Mg R^2}{\tau r } \right) } } }
\]



\solutionhead{3} \textbf{Potential energy of gas in gravitational field.}  

I looked at the solution from Ed Groth \cite{EGroth1999} for Problem 3, Homework No. 4, Physics 301 for a hint, which prompted me to recall that diffusive equilibrium between 2 systems means their chemical potential $\mu_i$'s are equal.  

\footnote{John D. Naud via Ed Groth, Problem 3, Homework No. 4, Physics 301, Fall 1999, \url{http://grothserver.princeton.edu/~groth/phys301f99/hw04sol.pdf} }

Problem 3 \textbf{Potential energy of gas in gravitational field} of Chapter 5: Chemical Potential and Gibbs Distribution from Kittel and Kroemer (1980) \cite{CKittelHKroemer1980}: ``Consider a column of atoms each of mass $M$ at temperature $\tau$ in a uniform gravitational field $g$.''

Consider an infinitesimal layer at height $h$ and at height $h+dh$.  At each height $h$, the layers of atoms are in \emph{diffusive} and \emph{thermal equilibrium}, and so
\[
\tau(h+dh) = \tau(h) = \tau
\]
(thermal equilibrium) and 
\[
\mu(h+dh) = \mu(h)
\]
(diffusive equilibrium).  

Now for chemical potential $\mu = \mu(\tau,V,N)$,
\[
\mu = \mu_{\text{int}} + \mu_{\text{ext}} = \tau \ln{ \left( n \left( \frac{ 2\pi \hbar^2 }{M\tau} \right)^{3/2} \right) } + Mgh
\]

Thus, to first order,
\[
\tau \frac{1}{n(h) } \frac{ dn(h) }{dh} = -Mg \text{ or } n(h) = n(0) \exp{ \left( \frac{-Mgh}{\tau} \right) }
\]

The thermal average potential energy per atom is 
\[
\begin{gathered}
  \langle U \rangle = \frac{ \int_0^{\infty} dh Mgh n(h) }{ \int_0^{\infty} dh n(h) } = Mg \frac{ \left. \left[ \frac{ h \exp{ \left( \frac{-Mgh}{\tau} \right) } }{ -Mg/\tau } + - \frac{ \exp{ \left( \frac{-Mgh}{\tau} \right) } }{ (-Mg/\tau)^2 } \right] \right|_0^{\infty} }{ \left. \left( \frac{ \exp{ \left( \frac{-Mgh}{\tau } \right) } }{ -Mg/\tau} \right) \right|_0^{\infty} } = Mg \frac{\tau}{Mg} = \boxed{ \tau = \langle U \rangle }
\end{gathered}
\]

From Ch.3, Eq. (64) of Kittel and Kroemer (1980) \cite{CKittelHKroemer1980}, 
\[
U = \frac{ \sum_{\mathbf{n}} \epsilon_{\mathbf{n}} \exp{ \left[ -\epsilon_{\mathbf{n}} /\tau \right] } }{Z_1} = \tau^2 \frac{ \partial \ln{Z_1}}{\partial \tau } = \frac{3\tau}{2}
\]
So the total heat capacity per atom $C$ is 
\[
C = \frac{ \partial E}{ \partial \tau} = \frac{ \partial ( \langle U \rangle + U ) }{ \partial \tau } = \frac{ \partial }{ \partial \tau} ( \tau + \frac{3\tau}{2} ) = \boxed{ \frac{5}{2}  = C}
\]

\solutionhead{4} $\mu = \mu_{\text{int}} + \mu_{\text{ext}} = \tau \ln{ \left( \frac{n}{n_Q} \right) } + \mu_{\text{ext}}$

\[
\mu_{\text{ext}} = \tau \ln{ \left( \frac{10^4 n }{n_Q} \right) } - \tau \ln{ \left( \frac{n}{n_Q} \right) } = \tau \ln{ 10^4} = 300 \, K * 0.0000861733 * \, eV/K \ln{ 10^4} = 0.24 \, eV
\]

Note that I checked this solution afterwards with the one by John D. Naud, via Ed Groth's materials \cite{EGroth1999}, for Problem 4, Homework No. 4, Physics 301 \footnote{John D. Naud via Ed Groth, Problem 4, Homework No. 4, Physics 301, Fall 1999, \url{http://grothserver.princeton.edu/~groth/phys301f99/hw04sol.pdf} }.  




\solutionhead{5} \textbf{Magnetic concentration}.  

This problem pertains to Figure 5.6.  Recall the example on pp. 127 of Kittel and Kroemer (1980) \cite{CKittelHKroemer1980}, \textbf{Example: Chemical potential of mobile magnetic particles in a magnetic field}.  ``Consider system of $N$ identical particles with magnetic moment $m$.''  

Now
\[
\begin{aligned}
  & \mu_{\text{tot}}(\uparrow) = \tau \ln{ \left( \frac{n_{\uparrow}}{n_Q} \right) } - mB = \tau \ln{n_{\uparrow}} + \frac{3\tau}{2} \ln{ \left( \frac{2\pi \hbar^2}{M\tau} \right) } - mB  \\
  & \mu_{\text{tot}}(\uparrow)(B=0) = \tau \ln{ \left( n_{\uparrow}(B=0) \right) } + \frac{3\tau}{2} \ln{\left( \frac{2\pi \hbar^2}{ M\tau }\right) }   \\
  & \mu_{\text{tot}}(\uparrow)(B=20 \, \text{kgauss}) = \tau \ln{ ( n_{\uparrow}(B=20\, \text{kgauss}))} + \frac{3\tau}{2} \ln{\left( \frac{2\pi \hbar^2}{M\tau} \right) } - m (20 \text{ \, kgauss} )
\end{aligned}
\]

When $\begin{aligned} & \quad \\
  & n_{\uparrow}(B=0) = 2\times 10^7 \, \text{cm}^{-3} \\
  & n_{\uparrow}(B=20 \, \text{ kgauss} )  = 2\times 10^9 \, \text{cm}^{-3} \end{aligned}$, 

\[
\mu_{\text{tot}}(\uparrow)(B=0) = \mu_{\text{tot}}(\uparrow)(B=20\, \text{ kgauss} ) 
\]
so
\[
\tau \ln{ \left( \frac{ 2\times 10^9 \, \text{cm}^{-3} }{ 2\times 10^7 \, \text{cm}^{-3} } \right)} = m(20 \, \text{kgauss} )
\]
Given $\tau =300 \, K$, 
\[
m = \frac{300 \, K  \cdot k_B \ln{ (10^2) } }{ 20 \, \text{kgauss} } = 9.54 \times 10^{-18} \, \text{erg}/\text{gauss} = 10^3 \text{ Bohr magnetons }
\]
where we take the Bohr magneton $\mu_B \equiv e\hbar/2mc$ to be $9.274 \times 10^{-21} \, \text{erg}/\text{gauss}$.  

\solutionhead{6} \textbf{Gibbs sum for a two level system}.  
Recall that 
\[
\mathcal{Z}(\mu, \tau) = \sum_{N=0}^{\infty} \sum_{s(N)} \exp{ [ (N \mu - \epsilon_{s(N)} )/\tau ]} = \sum_{ASN} \exp{ [(N\mu - \epsilon_{s(N)})/\tau ]}; \quad \quad \lambda = \exp{ \left( \frac{\mu}{\tau} \right)}
\]
\begin{itemize}
\item[(a)] $\mathcal{Z} = 1 + \lambda + \lambda \exp{ \left( \frac{- \epsilon}{\tau} \right)}$,  $1$ for $N=0$; for $N=1$, $\lambda + \lambda \exp{ \left( \frac{-\epsilon}{\tau} \right)}$.  $\lambda$ for $\epsilon =0$.  
\item[(b)] $\langle N \rangle = \frac{ 0 (1) + (1) (\lambda + \lambda \exp{ \left( \frac{-\epsilon}{\tau} \right)} ) }{ \mathcal{Z}} = \frac{ \lambda ( 1 + \exp{ (-\epsilon/\tau )} ) }{ \mathcal{Z}}$
\item[(c)] $\langle N(\epsilon) \rangle = \frac{ 0 (1) + (0) (\lambda) + \lambda \exp{ (-\epsilon/\tau) }  }{ \mathcal{Z}} = \frac{ \lambda \exp{ (-\epsilon/\tau) } }{ \mathcal{Z}}$
\item[(d)]$\langle \epsilon \rangle = \epsilon \langle N(\epsilon) \rangle = \frac{ \lambda \epsilon \exp{ (-\epsilon/\tau)} }{\mathcal{Z}}$  
\item[(e)] 
\[
\mathcal{Z} = 1 + \lambda + \lambda \exp{ (-\epsilon/\tau) } + \lambda^2 \exp{ (-\epsilon/\tau) } = (1+ \lambda)(1 + \lambda \exp{ ( -\epsilon/\tau) }
\]
where we considered the possibility the orbital at $0$ and at $\epsilon$ are each occupied by one particle at the same time.  \\
So that, for total energy being $\epsilon$, $\exp{ [ (2\mu - \epsilon)/\tau ]}=   (\exp{ \left( \frac{\mu}{\tau} \right)})^2 e^{-\epsilon/\tau} = \lambda^2 e^{-\epsilon/\tau}$.  
\end{itemize}

\solutionhead{7} \textbf{States of positive and negative ionization}.  
\[
\begin{aligned}
  & \mathcal{Z} = e^{ \frac{ \delta}{2\tau} } + \lambda e^{\frac{ \Delta }{ 2 \tau }} + \lambda e^{\frac{ -\Delta }{ 2 \tau }}  + \lambda^2 e^{ \frac{ -\delta}{2\tau} } \\
  &  \partial_{\lambda} \mathcal{Z} = 0 + e^{ \Delta/ 2\tau } + e^{\frac{ -\Delta }{ 2 \tau }} + 2 \lambda e^{\frac{ -\delta }{ 2 \tau }} \\ 
  & \lambda \partial_{ \partial} \ln{ \mathcal{Z}} = \frac{ \lambda ( e^{ \Delta/ 2\tau} + e^{-\Delta/2\tau} + 2 \lambda e^{-\delta/ 2\tau} ) }{ e^{ \frac{ \delta}{2\tau} } + \lambda e^{\frac{ \Delta }{ 2 \tau }} + \lambda e^{\frac{ -\Delta }{ 2 \tau }}  + \lambda^2 e^{ \frac{ -\delta}{2\tau} } } = 1 
\end{aligned}
\]
\[
\lambda e^{\Delta/ 2\tau} + \lambda e^{-\Delta/2\tau} + 2\lambda^2 e^{-\delta/2\tau} = e^{\delta/2\tau} = \lambda e^{\Delta/2\tau} + \lambda e^{-\Delta/2\tau} + \lambda^2 e^{-\delta/2\tau} \quad \quad \Longrightarrow \begin{gathered} \lambda^2 e^{-\delta/2\tau} = e^{\delta/2\tau} \\ \lambda^2 = e^{\delta/\tau} \text{ or } 2 \ln{\lambda } = \delta/\tau \end{gathered} 
\]
\[
\boxed{ 2 \tau \ln{\lambda } = \delta }
\]

\solutionhead{8} \textbf{Carbon monoxide poisoning}.  \begin{itemize}
\item[(a)]$\mathcal{Z} = 1 +  \lambda(O_2) e^{-\epsilon_A/\tau}$
\[
P(O_2) =  \frac{ \lambda(O_2) e^{-\epsilon_A}{\tau} }{ 1 + \lambda(O_2) e^{-\epsilon_A/\tau} } = 0.9 \text{ or } 0.9 = 0.1 \lambda(O_2) \exp{ \left( \frac{ - \epsilon_A}{\tau} \right) }
\]

\[
\Longrightarrow \ln{ \left( \frac{ 9 }{ \lambda(O_2) } \right)} = \frac{- \epsilon_A}{\tau}; \text{ or } \tau \ln{ \left( \frac{ \lambda(O_2)}{ 9} \right) } = \epsilon_A 
\]
\[
\boxed{ \epsilon_A = (k_B T) \ln{ \left( \frac{ \lambda(O_2)}{ 9 } \right)} } = (8.617 \times 10^{-5} \, \frac{eV}{k} )(273 + 39) \ln{ \left( \frac{10^{-5}}{9} \right)} = - 0.3686 \, eV
\]
\item[(b)] $\mathcal{Z} = 1 + \lambda(O_2) e^{- \epsilon_A/\tau} + \lambda(CO) e^{-\epsilon_B/\tau}$

\[
P(O_2) = \frac{ \lambda(O_2) e^{-\epsilon_A/\tau} }{ 1 + \lambda(O_2) e^{- \epsilon_A/\tau} + \lambda(CO) e^{-\epsilon_B /\tau} } = 0.1 \Longrightarrow 0.1 + 0.1 \lambda(CO) e^{-\epsilon_B/\tau} = 0.9 \lambda(O_2) e^{-\epsilon_A/\tau}
\]
\[
\ln{ \left( \frac{ 9 \lambda(O_2) e^{-\epsilon_A/\tau} -1}{ \lambda(CO) } \right)} = \frac{ -\epsilon_B}{\tau} \text{ or } \boxed{ \epsilon_B = \tau \ln{ \left( \frac{ \lambda(CO) }{ 9 \lambda(O_2) e^{-\epsilon_A/\tau} -1 } \right) } } = -0.5511 \, eV
\]
\end{itemize}

\solutionhead{9} \textbf{Absorption of $O_2$ in a magnetic field}.  Recall that $2j + 1 = $ total number of spin states.  \\
$2(1) + 1= 3$.  \\
Now $U = - \mathbf{m}\cdot \mathbf{B}$.  

\[
\begin{aligned}
  & \mathcal{Z} = 1 + \lambda(O_2) e^{ \frac{ - \epsilon + \mu_B B}{ \tau} } + \lambda(O_2) e^{ -\frac{ \epsilon_A}{\tau} } + \lambda(O_2) e^{ \frac{ - (\epsilon_A + \mu_B B) }{ \tau} } = \\
  & \mathcal{Z} = 1 + \lambda(O_2) e^{- \epsilon_A /\tau} ( 1 + 2 \cosh{ \left( \frac{ \mu_B B}{ \tau} \right) } )
\end{aligned}
\]

\[
\begin{gathered}
  0.91 = \frac{ \lambda(O_2) e^{- \epsilon_A/\tau} ( 1 + 2 \cosh{ \left( \frac{ \mu_B B}{ \tau } \right) } ) }{ 1 + \lambda(O_2) e^{ - \epsilon_A/\tau} ( 1 + 2 \cosh{ \left( \frac{ \mu_B B }{ \tau } \right) } ) } \text{ or } 0.91 = 0.09 \lambda(O_2) e^{ - \epsilon_A/\tau} ( 1 + 2 \cosh{ \left( \frac{ \mu_B B}{ \tau} \right) } ) \\ 
  \frac{1}{2} \left( \frac{ 91}{ 9} \frac{1}{ \lambda(O_2) } e^{ \epsilon_A/\tau} -1 \right) = \cosh{ \left( \frac{ \mu_B B}{\tau} \right) } \text{ or } B = \frac{ \tau}{ \mu_B} arccosh{ \left( \frac{1}{2} \left( \frac{91}{9} \frac{1}{ \lambda(O_2) } e^{ \epsilon_A/\tau} -1 \right) \right) } 
\end{gathered}
\]

The Gibbs sum in the limit of zero magnetic field will differ from that of Problem 8 because there the spin multiplicity of the bound state was neglected.  

\[
\begin{gathered}
  \mathcal{Z} = 1 + 3 \lambda(O_2) e^{- \epsilon_A/\tau} \\ 
  P(O_2) = \frac{ 3 \lambda(O_2) e^{-\epsilon_A/\tau} }{ 1 + 3 \lambda(O_2) e^{- \epsilon_A/\tau}} = 0.9 \text{ or } 0.9 = 0.3 \lambda(O_2) e^{ -\epsilon_A/\tau} \\
   \ln{ \frac{3}{\lambda(O_2)} } =  - \epsilon_A/\tau  \text{ or } \epsilon_A = \tau \ln{ \left( \frac{ \lambda(O_2) }{ 3} \right) } = -0.6227
\end{gathered}
\]
for $T = 300 \, K$, so that $\tau = 0.049375 \, eV$.  

\[
\Longrightarrow B = \frac{ \tau}{ \mu_B} 0.59927 = \frac{ 0.049375 \cdot 0.59927 }{ 5.7884 \times 10^{-11} \times 10^6 \, eV/T } = 511.1 \, T
\]





\solutionhead{10} \textbf{Concentration fluctuations}.  
\begin{itemize}
\item[(a)]
Recall that $\mathcal{Z} = \sum_{N=0}^{\infty} \sum_{s(N)} \exp{ [ (N \mu - \epsilon_{s(N)} )/\tau ] }$.  
\[
\frac{ \partial^2 \mathcal{Z} }{ \partial \mu^2}   = \sum_{ASN} \left( \frac{N}{\tau} \right)^2 \exp{ [ (N \mu - \epsilon_{s(N)} )/\tau ] } \quad \quad \quad \, \langle N^2 \rangle = \frac{ \tau^2 }{\mathcal{Z}} \frac{ \partial^2 \mathcal{Z}}{ \partial \mu^2 } 
\]
\item[(b)] 
\[
\frac{ \partial }{ \partial \mu} \langle N \rangle = \frac{ \partial }{ \partial \mu} \left( \frac{ \tau}{ \mathcal{Z} } \left( \frac{ \partial \mathcal{Z}}{ \partial \mu} \right)_{ \tau, V} \right) =\tau \left( \frac{ -  \left( \frac{ \partial \mathcal{Z}}{ \partial \mu} \right)^2_{\tau, V} }{ \mathcal{Z}^2} + \frac{ \left( \frac{ \partial^2 \mathcal{Z}}{  \partial \mu^2} \right)_{\tau, V} }{ \mathcal{Z}} \right) 
\]
\[
\Longrightarrow \tau \frac{ \partial}{ \partial \mu} \langle N \rangle = \langle N^2 \rangle - \langle N \rangle^2 = \langle ( \Delta N)^2 \rangle
\]
\end{itemize}

\solutionhead{11} \textbf{ Equivalent definition of chemical potential}.  

Recall that 
\[
d\sigma = \left( \frac{ \partial \sigma}{ \partial U} \right)_{V,N} dU + \left( \frac{ \partial \sigma}{ \partial V} \right)_{U,N} dV + \left( \frac{ \partial \sigma}{ \partial N } \right)_{U,V} dN \quad \quad \quad \, (31)
\]
\[
\mu = - \tau \left( \frac{ \partial \sigma}{ \partial N } \right)_{U,V} \quad \quad \quad \, (35)
\]
Consider when $d\sigma = dV =0$.  
\[
\begin{gathered}
  \Longrightarrow \frac{ - \mu }{ \tau} dN + \left( \frac{ \partial \sigma}{ \partial U } \right)_{V,N} dU = 0 \\
  \left( \frac{ \partial \sigma}{ \partial U} \right)_{V,N} \frac{dU}{dN} = \frac{ \mu}{ \tau}
\end{gathered}
\]

Note that $\frac{ dU}{ dN} = \left( \frac{ \partial U}{ \partial N} \right)_{\sigma, V}$.  

Using the definition, $\left( \frac{ \partial \sigma}{ \partial U} \right)_{N,V} \equiv \frac{1}{\tau}$, so then
\[
\boxed{ \mu = \left( \frac{ \partial U}{ \partial N} \right)_{\sigma, V}}
\]

Now $F = U - \tau \sigma$, by definition.  Consider the thermodynamic identity, $dU = \tau d\sigma - p dV + \mu dN$.  

\[
\begin{gathered}
\begin{aligned}
  dF & = dU - \sigma d\tau - \tau d\sigma = \tau d\sigma - pdV + \mu dN - \sigma d\tau - \tau d\sigma = \\
  & = - pdV + \mu dN - \sigma d\tau
\end{aligned} \\
\Longrightarrow \left( \frac{ \partial F}{ \partial N} \right)_{ \tau, V} = \mu(\tau, V, N)
\end{gathered}
\]

Likewise, $\mu = \left( \frac{   \partial U}{ \partial N} \right)_{\sigma, V}$ is equivalent to $\mu(\tau, V, N) = \left( \frac{ \partial F}{ \partial N} \right)_{\tau, V}$ through the thermodynamic identity as well.  
\[
dU = \tau d\sigma - p dV + \mu dN
\]
$\sigma, V$ constant.  $d\sigma, dV = 0$.  
\[
\Longrightarrow \left( \frac{ \partial U}{\partial N} \right)_{\sigma, V} = \mu 
\]

\solutionhead{12} \textbf{Ascent of sap in trees}.  Given: relative humidity $r = 0.9$.  $T = 25^{\circ} \, C$ \\

$n_0 = $ concentration in saturated air that stands immediately above pool of water of water vapor in air.  \\
$r n_0 = $ actual concentration of water vapor in air at uppermost leaves is $rn_0$.   \\

At pool, $\mu_{H_2O} = \mu_{\text{vapor}}$ for diffusion equilibrium.   \\

Same condition at uppermost leaves, otherwise there's evaporation: \\
$\mu_{\text{sap}} = \mu_{\text{vapor}}(h)$.  

$\mu_{\text{sap}} = \mu_{H_2O}$ (no flow going on in water).  Thus, $\mu_{\text{vapor}}(h) = \mu_{\text{vapor}}(0)$ and treat water vapor as an ideal gas.  

\[
\begin{gathered}
  \tau \ln{ \left( \frac{ n(h)}{ n_Q} \right) } + mgh = \tau \ln{ \left( \frac{n(0)}{n_Q } \right)} \\ 
  \tau \ln{ \left( \frac{ r n_0}{n_0} \right) } = - mgh \Longrightarrow \tau \ln{ \left( \frac{1}{r} \right) } = mgh \text{ or } h = \frac{ \tau}{mg} \ln{ \left( \frac{1}{r} \right) }
\end{gathered}
\]

\[
h  = \frac{ ( 1.381 \times 10^{-23} \, \frac{1}{K} )(298 \, K)( kg\cdot \frac{m^2}{s^2} ) \ln{ \left( \frac{10}{9} \right) } }{ (18 \, amu )\left( \frac{1.67 \times 10^{-27} \, kg}{ 1 \, amu} \right)(9.8 \, m/s^2) } = 147.1 \, m 
\]

\solutionhead{13} \textbf{Isentropic expansion.}  
\begin{itemize}
\item[(a)]
\item[(b)] Recall that $F = - \tau [ N \ln{Z_1} - \ln{N!} ]$ where $Z_1 = n_Q V = \left( \frac{ M \tau}{ 2\pi \hbar^2 } \right)^{3/2} V$ and \\
$\sigma = -\left( \frac{ \partial F}{ \partial \tau} \right)_{V,W}$.  
\[
\begin{aligned}
  \frac{ \partial F}{ \partial \tau} & = - [ N \ln{Z_1} - \ln{N!} ] - \tau [ N \frac{1}{Z_1} \left( \frac{ M }{ 2\pi \hbar^2 } \right)^{3/2} \frac{3}{2} \tau^{1/2} V ] = \\
  & = - [  N \ln{Z_1} - \ln{N!} ] - [ N \frac{1}{Z_1} n_Q \frac{3}{2} V ] = - [ N \ln{Z_1} - \ln{N!} ] - \frac{3}{2} N
\end{aligned}
\]
With $N$ constant and $Z_1 = \left( \frac{M}{ 2\pi \hbar^2 } \right)^{3/2} \tau^{3/2} V$, then for an isentropic expansion, $\tau V^{2/3}$ must remain constant.  

\end{itemize}

\solutionhead{14} \textbf{Multiple binding of $O_2$}.  

\begin{itemize}
\item[(a)] Be wary of the multiplicity, how you count, each of the energy states.  
\[
\mathcal{Z} = 1 + 4 \lambda e^{-\epsilon/\tau} + \binom{4}{2} \lambda^2 e^{-2\epsilon/\tau} + \binom{4}{3} \lambda^3 e^{-3\epsilon/\tau} + \lambda^4 e^{-4\epsilon/\tau} = (1 + \lambda e^{-\epsilon/\tau})^4
\]
\[
\boxed{ P(\epsilon ) = \frac{ 4 \lambda e^{-\epsilon/\tau}}{ (1 + \lambda e^{-\epsilon/\tau})^4} }
\]
\item[(b)] \[
P(4\epsilon) = \frac{ \lambda^4 e^{-\epsilon/\tau}}{ (1+ \lambda e^{-\epsilon/ \tau})^4} = \frac{ e^{-\epsilon /\tau}}{ \left( \frac{1}{\lambda + e^{-\epsilon/ \tau} } \right)^4}
\]
\end{itemize}






